\par\SourceLeaf{425}%
\noindent\hspace*{1.2em}We suppose the points are ordered so that they run around the triangle they enclose in the same sense as the points $(1),(2),(3)$. Let $\lambda$ further be the pole of the great circle $L'L''$ that lies on the same side as $L$. From our preceding lemma we then have
\[
\begin{aligned}
y'z''-z'y''&=\sin L'L''.\cos\lambda(1)\\
z'x''-x'z''&=\sin L'L''.\cos\lambda(2)\\
x'y''-y'x''&=\sin L'L''.\cos\lambda(3).
\end{aligned}
\]
Multiplying these equations by $x,y,z$ and adding the products gives
\[
xy'z''+x'y''z+x''yz'-xy''z'-x'yz''-x''y'z=\sin L'L''.\cos\lambda L,
\]
for which, by the familiar rules of spherical trigonometry, we can also write
\[
\begin{aligned}
&\sin L'L''.\sin LL''.\sin L'\\
&=\sin L'L''.\sin LL'.\sin L''\\
&=\sin L'L''.\sin L'L''.\sin L
\end{aligned}
\]
if $L,L',L''$ denote the three angles of the spherical triangle $LL'L''$. One can also easily see that this quantity is one sixth of the pyramid formed between the center of the sphere and the points $L,L',L''$ (and is \textls[100]{positive}, if etc.).

\Needspace{5\baselineskip}
\begin{center}8.\end{center}
\noindent\hspace*{1.2em}The nature of a curved surface is expressed by an equation between the coordinates of its points, which we write as
\[
f(x,y,z)=0
\]
Let the total differential of $f(x,y,z)$ be
\[
=Pdx+Qdy+Rdz,
\]
where $P,Q,R$ will be functions of $x,y,z$. We will always distinguish two sides of the surface, which we will call the upper and the lower side. Generally speaking, as one passes through\space%
